\documentclass[aspectratio=169]{beamer}

\usepackage{algorithm}
\usepackage{algpseudocode}
\usepackage{tikz}
\usetikzlibrary{arrows.meta,positioning,shapes.geometric}
\usepackage{calc}
\usepackage{polyglossia}
\setmainlanguage{vietnamese}
\usepackage{amsthm,amsmath,amssymb}
\usepackage{graphicx}
\usepackage{booktabs}
\usepackage{tabularx}
\usepackage{array}
\usepackage{subcaption}
\usepackage{xcolor}

\newtheorem{thm}{Định lý}
\newtheorem{defn}{Định nghĩa}
\renewcommand{\proofname}{\textbf{Chứng minh}}
\renewcommand{\qedsymbol}{$\blacksquare$}
\setbeamertemplate{theorems}[numbered]

%% TYPESET
\makeatletter
\newcommand{\leqnomode}{\tagsleft@true}
\newcommand{\reqnomode}{\tagsleft@false}
\makeatother
\let\Mathcal\mathcal
\IfFileExists{dutchcal.sty}{\usepackage{dutchcal}}{}
\newcommand{\norm}[1]{\left\lVert#1\right\rVert}
\newcommand{\abs}[1]{\left\lvert#1\right\rvert}
\newcommand{\gr}{\operatorname{gr}}
\newcommand{\rank}{\operatorname{rank}}
\newcommand{\toto}{\rightrightarrows}
\newcommand{\F}{\Mathcal{F}}
\newcommand{\K}{\mathbb{K}}
\newcommand{\R}{\mathbb{R}}
\def\C{\mathbb{C}}
\newcommand{\dom}{\operatorname{dom}}
\newcommand{\im}{\operatorname{im}}
\newcommand{\g}{\mathcal{g}}
\global\long\def\MP{\text{mp}}

% Dùng theme HUST nếu các tệp theme có trong project Overleaf.
% Fallback Madrid giúp slide vẫn biên dịch được khi chỉ upload file .tex.
\IfFileExists{beamerthemehust.sty}{
    \usetheme{hust}
}{
    \usetheme{Madrid}
    \definecolor{hustred}{RGB}{177,18,38}
    \setbeamercolor{structure}{fg=hustred}
    \setbeamercolor{frametitle}{fg=white,bg=hustred}
}

\setbeamertemplate{caption}[numbered]
\setbeamertemplate{navigation symbols}{}
\setbeamertemplate{bibliography item}[text]

% Giống báo cáo: ảnh có thể nằm trong overleaf_assets/ ở repository
% hoặc được upload trực tiếp vào thư mục gốc của Overleaf.
\graphicspath{{overleaf_assets/}{./}}
\newcommand{\safeincludegraphics}[2][]{%
    \IfFileExists{#2}{%
        \includegraphics[#1]{#2}%
    }{%
        \IfFileExists{overleaf_assets/#2}{%
            \includegraphics[#1]{overleaf_assets/#2}%
        }{%
            \fbox{\parbox[c][3.2cm][c]{0.86\linewidth}{%
                \centering\textbf{Thiếu tệp hình}\\[0.15cm]
                \texttt{\detokenize{#2}}\\[0.15cm]
                Tải tệp từ \texttt{overleaf\_assets} lên Overleaf.%
            }}%
        }%
    }%
}

\newcommand{\Image}[1]{%
    \sbox0{\safeincludegraphics[height=0.65\paperheight]{#1}}%
    \ifdim\wd0 < \textwidth
        \safeincludegraphics[height=0.65\paperheight]{#1}%
    \else
        \safeincludegraphics[width=\textwidth]{#1}%
    \fi%
}

\newcommand{\modelname}{BiLSTM + Attention}
\newcommand{\codeword}[1]{\texttt{\detokenize{#1}}}
\newcommand{\tokenbox}[2]{%
    \tikz[baseline=(tok.base)]{
        \node[
            rounded corners=2pt,
            fill=#1,
            inner xsep=4pt,
            inner ysep=2pt
        ] (tok) {\strut #2};
    }%
}

\title{\textbf{PHÂN LOẠI CẢM XÚC ĐA MỨC\\TRÊN AMAZON REVIEWS}}
\subtitle{Mô hình BiLSTM kết hợp Additive Attention}
\author{Nhóm 5}
\institute{Trường Công nghệ Thông tin và Truyền thông\\Đại học Bách khoa Hà Nội}
\date{2026}

\begin{document}

\begin{frame}[plain,noframenumbering]
    \maketitle
\end{frame}

\begin{frame}{\textbf{\large Thành viên nhóm}}
    \vspace*{-0.45cm}
    \begin{large}
        \begin{enumerate}
            \item GVHD: PGS. TS. Ngô Văn Linh
            \item Trần Lê Ngọc Tâm
            \hfill \makebox[4.2cm][l]{MSSV: 202416346}
            \item Nguyễn Phong
            \hfill \makebox[4.2cm][l]{MSSV: 202400066}
            \item Nguyễn Ngọc Tuấn Anh
            \hfill \makebox[4.2cm][l]{MSSV: 202400029}
        \end{enumerate}
    \end{large}
\end{frame}

\begin{frame}{\textbf{\large Nội dung}}
    \tableofcontents
\end{frame}

% =========================================================
\section{1. Bài toán và dữ liệu}

\begin{frame}{\textbf{\large Bài toán và mục tiêu}}
    \begin{columns}[T,onlytextwidth]
        \begin{column}{0.55\textwidth}
            \begin{block}{Bài toán}
                Phân loại mỗi review vào \textbf{một trong 5 lớp}:
                \[
                    y\in\{1,2,3,4,5\}\text{ sao}.
                \]
                Đây là \textbf{multiclass classification}, không phải hồi quy.
            \end{block}

            \begin{itemize}
                \item Đầu ra: Softmax 5 lớp.
                \item Loss: sparse categorical cross-entropy.
                \item Dự đoán: lớp có xác suất lớn nhất.
            \end{itemize}
        \end{column}

        \begin{column}{0.41\textwidth}
            \begin{alertblock}{Mục tiêu chính}
                \begin{itemize}
                    \item Nắm bắt ngữ cảnh hai chiều.
                    \item Tập trung vào token quan trọng.
                    \item So sánh với SimpleRNN và Transformer Encoder.
                    \item Đánh giá cả lớp thiểu số, không chỉ Accuracy.
                \end{itemize}
            \end{alertblock}
        \end{column}
    \end{columns}
\end{frame}

\begin{frame}{\textbf{\large Dataset: Amazon Reviews}}
    \begin{columns}[T,onlytextwidth]
        \begin{column}{0.50\textwidth}
            \begin{block}{Nguồn dữ liệu}
                \textbf{Amazon Reviews Dataset}\\[0.12cm]
                \href{https://www.kaggle.com/datasets/dongrelaxman/amazon-reviews-dataset}
                {\textcolor{blue}{Mở trang dữ liệu trên Kaggle}}
            \end{block}

            \begin{itemize}
                \item Hai trường sử dụng: \codeword{Review Text}, \codeword{Rating}.
                \item Sau làm sạch: \textbf{20.405 review hợp lệ}.
                \item Chia phân tầng: 80\% train, 20\% validation.
                \item Validation: \textbf{4.081 mẫu}.
            \end{itemize}

            \begin{alertblock}{Baseline dễ gây hiểu nhầm}
                Nếu luôn đoán \textbf{1 sao}, mô hình đã đạt
                \textbf{63,49\% Accuracy} mà chưa học được năm mức cảm xúc.
            \end{alertblock}
        \end{column}

        \begin{column}{0.46\textwidth}
            \centering
            \small
            \begin{tabular}{lrr}
                \toprule
                \textbf{Lớp} & \textbf{Số mẫu} & \textbf{Tỷ lệ} \\
                \midrule
                1 sao & 2.591 & 63,49\% \\
                2 sao & 241 & 5,91\% \\
                3 sao & 165 & 4,04\% \\
                4 sao & 239 & 5,86\% \\
                5 sao & 845 & 20,71\% \\
                \bottomrule
            \end{tabular}

            \vspace{0.25cm}
            \begin{alertblock}{Lớp thiểu số}
                Ba lớp 2--4 sao chỉ chiếm \textbf{15,80\%}. Vì vậy kết quả phải
                đọc cùng \textbf{Macro F1}, \textbf{rating MAE} và F1 từng lớp.
            \end{alertblock}
        \end{column}
    \end{columns}
\end{frame}

\begin{frame}{\textbf{\large Vì sao cần Macro F1 và rating MAE?}}
    \begin{columns}[T,onlytextwidth]
        \begin{column}{0.48\textwidth}
            \begin{block}{Macro F1: công bằng giữa các lớp}
                \begin{align*}
                    F1_k
                    &=\frac{2P_kR_k}{P_k+R_k},\\
                    \mathrm{MacroF1}
                    &=\frac{1}{5}\sum_{k=1}^{5}F1_k.
                \end{align*}
                Mỗi lớp đóng góp đúng \textbf{20\%}, bất kể có 2.591 hay chỉ
                165 mẫu. Nếu bỏ qua lớp 2--4 sao, Macro F1 giảm mạnh.
            \end{block}
        \end{column}

        \begin{column}{0.48\textwidth}
            \begin{block}{MAE: đo độ xa của lỗi rating}
                \[
                    \mathrm{MAE}_{rating}
                    =\frac{1}{N}\sum_{i=1}^{N}|y_i-\hat y_i|.
                \]
                \centering
                $1\!\rightarrow\!2$: lỗi 1
                \qquad
                $1\!\rightarrow\!5$: lỗi 4
            \end{block}
        \end{column}
    \end{columns}

    \vspace{0.25cm}
    \begin{exampleblock}{Ba góc nhìn bổ sung cho nhau}
        \centering
        Accuracy: đúng hay sai
        \quad $\big|$ \quad
        Macro F1: có bỏ quên lớp thiểu số không
        \quad $\big|$ \quad
        MAE: sai xa bao nhiêu sao
    \end{exampleblock}
\end{frame}

\begin{frame}{\textbf{\large Pipeline xử lý dữ liệu}}
    \vspace*{-0.2cm}
    \begin{tikzpicture}[
        node distance=0.34cm,
        pstep/.style={
            draw=hustred,
            rounded corners,
            fill=hustred!7,
            minimum height=0.82cm,
            text width=2.48cm,
            align=center,
            font=\scriptsize
        },
        arrow/.style={-{Latex[length=2mm]},thick,hustred}
    ]
        \node[pstep] (clean) {\textbf{1. Làm sạch}\\lowercase, bỏ HTML, trùng/rỗng};
        \node[pstep,right=of clean] (split) {\textbf{2. Stratified split}\\80\% train -- 20\% validation};
        \node[pstep,right=of split] (fit) {\textbf{3. Fit tokenizer}\\chỉ trên train, có \codeword{<OOV>}};
        \node[pstep,right=of fit] (pad) {\textbf{4. Mã hóa}\\post-pad/truncate về 200 token};
        \draw[arrow] (clean) -- (split);
        \draw[arrow] (split) -- (fit);
        \draw[arrow] (fit) -- (pad);
    \end{tikzpicture}

    \vspace{0.28cm}
    \begin{exampleblock}{Một review đi qua pipeline}
        \small
        \begin{tabularx}{\linewidth}{>{\bfseries}p{2.1cm} X}
            Dữ liệu thô &
            \codeword{<b>Battery</b> is NOT good!!!} \quad Rating: 2 \\[0.12cm]
            Sau làm sạch &
            \tokenbox{blue!12}{battery}\ 
            \tokenbox{blue!12}{is}\ 
            \tokenbox{orange!22}{not}\ 
            \tokenbox{red!18}{good} \\[0.12cm]
            Token IDs &
            $[842,\ 17,\ 39,\ 64]$ \scriptsize (minh họa) \\[0.12cm]
            Sau padding &
            $[842,\ 17,\ 39,\ 64,\ 0,\ 0,\ldots,0]\in\mathbb{N}^{200}$ \\[0.12cm]
            Embedding &
            $X\in\mathbb{R}^{200\times100}$; mỗi ID được ánh xạ sang vector
            GloVe 100 chiều.
        \end{tabularx}
    \end{exampleblock}

    \vspace{0.12cm}
    \centering
    \footnotesize
    \textbf{Chống leakage:}
    chia dữ liệu trước, rồi mới học vocabulary từ train.
\end{frame}

\begin{frame}{\textbf{\large Tại sao chọn GloVe 6B, 100 chiều?}}
    \begin{columns}[T,onlytextwidth]
        \begin{column}{0.49\textwidth}
            \begin{block}{GloVe cung cấp gì?}
                \begin{itemize}
                    \item Học từ thống kê đồng xuất trên \textbf{6 tỷ token}
                    của Wikipedia 2014 + Gigaword 5.
                    \item Khoảng \textbf{400.000 từ}; có sẵn các bản
                    50/100/200/300 chiều.
                    \item Các từ gần nghĩa có vector gần nhau ngay trước khi
                    train trên Amazon Reviews.
                \end{itemize}
            \end{block}
        \end{column}

        \begin{column}{0.47\textwidth}
            \begin{alertblock}{Phù hợp với thí nghiệm này}
                \begin{itemize}
                    \item Tập train chỉ có \textbf{16.324 review}: học embedding
                    hoàn toàn từ đầu dễ thiếu dữ liệu.
                    \item \textbf{100d} cân bằng thông tin và chi phí:
                    $26.715\times100\approx2,67$ triệu tham số.
                    \item \codeword{trainable=True}: giữ tri thức pretrained,
                    sau đó thích nghi với ngôn ngữ review.
                \end{itemize}
            \end{alertblock}
        \end{column}
    \end{columns}

    \vspace{0.18cm}
    \begin{exampleblock}{Trade-off}
        \centering
        Nhanh và nhẹ hơn 300d, giàu thông tin hơn 50d; nhưng không phải
        embedding theo ngữ cảnh và có thể thiếu từ đặc thù sản phẩm.
    \end{exampleblock}
\end{frame}

% =========================================================
\section{2. BiLSTM kết hợp Attention}

\begin{frame}{\textbf{\large Vì sao chọn BiLSTM + Attention?}}
    \begin{columns}[T,onlytextwidth]
        \begin{column}{0.31\textwidth}
            \begin{block}{LSTM}
                \centering
                Ghi nhớ phụ thuộc dài tốt hơn SimpleRNN nhờ cell state và các cổng.
            \end{block}
        \end{column}

        \begin{column}{0.31\textwidth}
            \begin{block}{Bidirectional}
                \centering
                Mỗi token nhận ngữ cảnh từ cả bên trái và bên phải.
            \end{block}
        \end{column}

        \begin{column}{0.31\textwidth}
            \begin{block}{Attention}
                \centering
                Gán trọng số lớn hơn cho token quan trọng đối với sentiment.
            \end{block}
        \end{column}
    \end{columns}

    \vspace{0.55cm}
    \begin{exampleblock}{Ví dụ}
        ``The screen is excellent, \textbf{but} the battery is
        \textbf{disappointing}.''\\
        Mô hình cần đọc hai chiều và xác định phần nào ảnh hưởng mạnh hơn tới
        rating cuối cùng.
    \end{exampleblock}
\end{frame}

\begin{frame}{\textbf{\large Bên trong một ô LSTM}}
    \begin{columns}[c,onlytextwidth]
        \begin{column}{0.47\textwidth}
            \centering
            \safeincludegraphics[
                width=\linewidth,
                height=0.61\paperheight,
                keepaspectratio
            ]{lstm.png}

            \vspace{0.08cm}
            \scriptsize
            Cell state $c_t$ là ``đường nhớ''; ba cổng điều khiển luồng thông tin.
        \end{column}

        \begin{column}{0.51\textwidth}
            \scriptsize
            \begin{align*}
                f_t &= \sigma(W_fx_t+U_fh_{t-1}+b_f),\\
                i_t &= \sigma(W_ix_t+U_ih_{t-1}+b_i),\\
                \widetilde c_t &= \tanh(W_cx_t+U_ch_{t-1}+b_c),\\
                c_t &= f_t\odot c_{t-1}+i_t\odot\widetilde c_t,\\
                o_t &= \sigma(W_ox_t+U_oh_{t-1}+b_o),\\
                h_t &= o_t\odot\tanh(c_t).
            \end{align*}

            \footnotesize
            $f_t$: quên thông tin cũ;\quad
            $i_t$: ghi thông tin mới;\quad
            $o_t$: chọn phần được xuất.

            \vspace{0.15cm}
            \begin{block}{Ý nghĩa}
                Phép cộng trên $c_t$ tạo đường truyền gradient ổn định hơn
                SimpleRNN, giúp giữ tín hiệu sentiment qua chuỗi dài.
            \end{block}
        \end{column}
    \end{columns}
\end{frame}

\begin{frame}{\textbf{\large Kiến trúc mô hình chính}}
    \centering
    \begin{tikzpicture}[
        node distance=0.55cm and 0.85cm,
        layer/.style={
            draw,
            rounded corners,
            minimum height=1.22cm,
            text width=3.05cm,
            inner xsep=4pt,
            align=center,
            font=\small
        },
        arrow/.style={-{Latex[length=2.2mm]},thick}
    ]
        \node[layer,fill=gray!10] (input) {Input\\$(B,200)$};
        \node[layer,fill=blue!10,right=of input] (embed)
            {GloVe Embedding\\$26.715\times100$};
        \node[layer,fill=green!12,right=of embed] (bilstm)
            {BiLSTM\\$128+128$ units};

        \node[layer,fill=orange!18,below=of bilstm] (attn)
            {Masked Additive\\Attention};
        \node[layer,fill=purple!12,left=of attn] (dense)
            {Dense 64\\ReLU + L2};
        \node[layer,fill=red!12,left=of dense] (out)
            {Softmax 5 lớp\\$\hat y=\arg\max p_k$};

        \draw[arrow] (input) -- (embed);
        \draw[arrow] (embed) -- node[above,font=\scriptsize]{Dropout} (bilstm);
        \draw[arrow] (bilstm) -- node[right,font=\scriptsize]{Dropout} (attn);
        \draw[arrow] (attn) -- (dense);
        \draw[arrow] (dense) -- node[above,font=\scriptsize]{Dropout} (out);
    \end{tikzpicture}

    \vspace{0.25cm}
    \begin{columns}[T,onlytextwidth]
        \begin{column}{0.49\textwidth}
            \begin{itemize}
                \item $h_t=[\overrightarrow h_t;\overleftarrow h_t]
                \in\mathbb{R}^{256}$.
                \item BiLSTM output: $(B,200,256)$.
                \item Attention context: $(B,256)$.
            \end{itemize}
        \end{column}
        \begin{column}{0.47\textwidth}
            \begin{itemize}
                \item Dropout: 0,5; L2: 0,01.
                \item Tổng tham số: 2.988.817.
            \end{itemize}
        \end{column}
    \end{columns}
\end{frame}

\begin{frame}{\textbf{\large Attention làm gì với chuỗi hidden state?}}
    \vspace*{-0.15cm}
    \centering
    \begin{tikzpicture}[
        token/.style={
            rounded corners,
            minimum width=1.65cm,
            minimum height=0.72cm,
            align=center,
            font=\scriptsize
        },
        arrow/.style={-{Latex[length=2mm]},thick}
    ]
        \node[token,fill=blue!10] (the) at (0,0) {the\\$\alpha_1=0{,}03$};
        \node[token,fill=green!18] (screen) at (1.9,0)
            {screen\\$\alpha_2=0{,}10$};
        \node[token,fill=green!35] (excellent) at (3.8,0)
            {excellent\\$\alpha_3=0{,}22$};
        \node[token,fill=orange!18] (but) at (5.7,0)
            {but\\$\alpha_4=0{,}08$};
        \node[token,fill=red!20] (battery) at (7.6,0)
            {battery\\$\alpha_5=0{,}17$};
        \node[token,fill=red!55] (bad) at (9.5,0)
            {disappointing\\$\alpha_6=0{,}40$};

        \node[draw,rounded corners,fill=purple!10,align=center,font=\small]
            (context) at (4.75,-1.65)
            {Context vector $c$\\biểu diễn toàn review};

        \foreach \x in {the,screen,excellent,but,battery,bad}
            \draw[arrow,gray] (\x.south) -- (context.north);
    \end{tikzpicture}

    \vspace{0.12cm}
    \begin{columns}[T,onlytextwidth]
        \begin{column}{0.48\textwidth}
            \begin{block}{Không dùng hidden state cuối}
                BiLSTM trả về
                $H=[h_1,\ldots,h_{200}]$, với $h_t\in\mathbb{R}^{256}$.
                Attention học mức đóng góp của \textbf{từng vị trí}.
            \end{block}
        \end{column}
        \begin{column}{0.48\textwidth}
            \begin{block}{Trực giác}
                Token mang sentiment mạnh nhận $\alpha_t$ lớn. Context là
                trung bình có trọng số, nên vẫn giữ thông tin ở xa cuối câu.
            \end{block}
        \end{column}
    \end{columns}

    \scriptsize\raggedright
    Các trọng số trên chỉ là ví dụ minh họa cơ chế, không phải output đã trích
    từ artifact hiện tại.
\end{frame}

\begin{frame}{\textbf{\large Additive Attention: từ score đến context}}
    \begin{columns}[T,onlytextwidth]
        \begin{column}{0.57\textwidth}
            \small
            Với $H\in\mathbb{R}^{B\times200\times256}$:
            \begin{align*}
                u_t &= \tanh(W_ah_t+b_a),\\
                e_t &= v_a^\top u_t,\\
                \alpha_t &= \frac{\exp(e_t)}
                {\sum_{j=1}^{T}\exp(e_j)},\\
                c &= \sum_{t=1}^{T}\alpha_t h_t.
            \end{align*}

            \begin{block}{Vai trò của từng hàm}
                \begin{itemize}
                    \item $\tanh$: tạo phép chấm điểm phi tuyến.
                    \item $v_a^\top$: nén vector thành một score vô hướng.
                    \item Softmax: chuẩn hóa $\alpha_t\ge0$ và
                    $\sum_t\alpha_t=1$.
                    \item Tổng có trọng số: $(B,200,256)\rightarrow(B,256)$.
                \end{itemize}
            \end{block}
        \end{column}

        \begin{column}{0.39\textwidth}
            \begin{alertblock}{Mask padding}
                Trước Softmax:
                \[
                    e_t\leftarrow
                    \begin{cases}
                        e_t,&m_t=1,\\
                        -10^9,&m_t=0.
                    \end{cases}
                \]
                Vì $\exp(-10^9)\approx0$, token 0 không được góp vào context.
            \end{alertblock}

            \begin{block}{Tham số học được}
                \begin{itemize}
                    \item $W_a\in\mathbb{R}^{256\times256}$
                    \item $b_a\in\mathbb{R}^{256}$
                    \item $v_a\in\mathbb{R}^{256}$
                \end{itemize}
                Tổng: \textbf{66.048 tham số}.
            \end{block}
        \end{column}
    \end{columns}
\end{frame}

\begin{frame}{\textbf{\large Hàm mất mát: Cross-Entropy + L2}}
    \begin{columns}[T,onlytextwidth]
        \begin{column}{0.56\textwidth}
            \begin{block}{1. Khớp nhãn bằng Cross-Entropy}
                Softmax đổi logits thành phân phối xác suất:
                \[
                    p_{i,k}=\frac{\exp(z_{i,k})}
                    {\sum_{j=0}^{4}\exp(z_{i,j})},
                    \qquad \sum_{k=0}^{4}p_{i,k}=1.
                \]
                Với nhãn nguyên $y_i\in\{0,\ldots,4\}$:
                \[
                    \mathcal L_{\mathrm{CE}}
                    =-\frac{1}{N}\sum_{i=1}^{N}\log p_{i,y_i}.
                \]
                Bản \textbf{sparse} nhận trực tiếp nhãn nguyên, không cần one-hot.
            \end{block}

            \begin{block}{2. Phạt độ phức tạp bằng L2}
                \[
                    \mathcal L
                    =\mathcal L_{\mathrm{CE}}
                    +\lambda\!\left(
                    \lVert W_{\mathrm{dense}}\rVert_2^2+
                    \lVert W_{\mathrm{out}}\rVert_2^2\right),
                    \quad\lambda=0{,}01.
                \]
            \end{block}
        \end{column}

        \begin{column}{0.40\textwidth}
            \begin{exampleblock}{Tại sao ``chồng'' hai thành phần?}
                \begin{itemize}
                    \item Softmax + CE: tạo xác suất rồi tối đa hóa xác suất của
                    lớp đúng; sai mà quá tự tin sẽ bị phạt mạnh.
                    \item L2: hạn chế trọng số quá lớn, giảm overfitting.
                    \item Một thành phần học dữ liệu; một thành phần kiểm soát
                    độ phức tạp.
                \end{itemize}
            \end{exampleblock}

            \begin{alertblock}{Nhược điểm trong bài toán này}
                \begin{itemize}
                    \item CE coi năm lớp độc lập: sai $1\!\to\!2$ và
                    $1\!\to\!5$ đều là ``sai lớp''.
                    \item Không tự cân bằng lớp; artifact đang
                    \textbf{tắt class weights}.
                    \item $\lambda$ quá lớn có thể gây underfitting.
                \end{itemize}
            \end{alertblock}
        \end{column}
    \end{columns}
\end{frame}

\begin{frame}{\textbf{\large Loss tối ưu khác metric đánh giá}}
    \begin{columns}[T,onlytextwidth]
        \begin{column}{0.47\textwidth}
            \begin{block}{Khi huấn luyện}
                \centering
                $\mathcal L_{\mathrm{CE}}+\mathcal L_{\mathrm{L2}}$

                \vspace{0.18cm}
                \raggedright
                Trơn, khả vi và cho gradient trực tiếp tới toàn bộ xác suất
                Softmax. Adam tối thiểu hóa loss này.
            \end{block}

            \begin{block}{Cấu hình}
                \small
                Adam $10^{-3}$; batch 128; tối đa 10 epoch; Dropout 0,5;
                EarlyStopping theo \codeword{val_loss}.
            \end{block}
        \end{column}

        \begin{column}{0.49\textwidth}
            \begin{alertblock}{Khi đánh giá}
                \begin{itemize}
                    \item \textbf{Macro F1}: đo chất lượng đồng đều trên năm lớp.
                    \item \textbf{Rating MAE}: đo khoảng cách giữa hai rating.
                    \item Cả hai đang tính sau $\arg\max$, nên không thuận tiện
                    để tối ưu trực tiếp bằng gradient.
                \end{itemize}
            \end{alertblock}

            \begin{exampleblock}{Hướng cải thiện}
                Class-weighted CE hoặc focal loss cho lớp thiểu số; ordinal loss
                hoặc CE kết hợp expected-rating loss để đưa thứ tự 1--5 sao vào
                quá trình học.
            \end{exampleblock}
        \end{column}
    \end{columns}
\end{frame}

\begin{frame}{\textbf{\large Kết quả riêng của BiLSTM + Attention}}
    \begin{columns}[c,onlytextwidth]
        \begin{column}{0.57\textwidth}
            \centering
            \safeincludegraphics[width=\linewidth,height=0.39\paperheight,keepaspectratio]
            {bilstm_attention_training_history.png}

            \vspace{0.05cm}
            \scriptsize Learning curve trên 10 epoch
        \end{column}

        \begin{column}{0.39\textwidth}
            \centering
            \safeincludegraphics[height=0.39\paperheight,keepaspectratio]
            {bilstm_attention_confusion_matrix.png}

            \vspace{0.05cm}
            \scriptsize Confusion matrix trên 4.081 mẫu validation
        \end{column}
    \end{columns}

    \vspace{0.35cm}
    \centering
    \begin{beamercolorbox}[
        wd=\textwidth,
        sep=0.18cm,
        center,
        rounded=true
    ]{block body}
        \textbf{Accuracy 81,38\% \qquad Macro F1 0,3727
        \qquad Rating MAE 0,3328}
    \end{beamercolorbox}
\end{frame}

% =========================================================
\section{3. So sánh mô hình}

\begin{frame}{\textbf{\large So sánh chất lượng tổng hợp}}
    \begin{columns}[T,onlytextwidth]
        \begin{column}{0.54\textwidth}
            \centering
            \safeincludegraphics[width=\linewidth]{quality_metrics.png}
        \end{column}

        \begin{column}{0.44\textwidth}
            \centering
            \scriptsize
            \begin{tabular}{lrrr}
                \toprule
                \textbf{Mô hình} & \textbf{Acc.} & \textbf{Macro F1} &
                \textbf{MAE} \\
                \midrule
                SimpleRNN & 0,7493 & 0,3029 & 0,6288 \\
                \textbf{BiLSTM + Attn.} &
                \textbf{0,8138} & \textbf{0,3727} & \textbf{0,3328} \\
                Transformer & 0,7988 & 0,3336 & 0,4009 \\
                \bottomrule
            \end{tabular}

            \vspace{0.45cm}
            \begin{block}{Nhận xét}
                \begin{itemize}
                    \item BiLSTM + Attention tốt nhất ở mọi metric chất lượng.
                    \item Hơn Transformer 1,49 điểm \% Accuracy.
                    \item Giảm MAE 47,08\% so với SimpleRNN.
                \end{itemize}
            \end{block}
        \end{column}
    \end{columns}
\end{frame}

\begin{frame}{\textbf{\large F1-score theo từng mức rating}}
    \begin{columns}[T,onlytextwidth]
        \begin{column}{0.61\textwidth}
            \centering
            \safeincludegraphics[width=\linewidth]{per_class_f1.png}
        \end{column}

        \begin{column}{0.36\textwidth}
            \begin{alertblock}{Phát hiện quan trọng}
                \begin{itemize}
                    \item Cả ba mô hình nhận diện tốt lớp 1 và 5 sao.
                    \item SimpleRNN và Transformer có F1 bằng 0 ở lớp 2--4.
                    \item BiLSTM + Attention có dự đoán đúng lớp giữa, nhưng F1
                    vẫn dưới 0,07.
                \end{itemize}
            \end{alertblock}

            \small
            Accuracy cao \textbf{không đồng nghĩa} bài toán 5 lớp đã được giải
            quyết tốt.
        \end{column}
    \end{columns}
\end{frame}

\begin{frame}{\textbf{\large So sánh chi phí tính toán}}
    \begin{columns}[T,onlytextwidth]
        \begin{column}{0.62\textwidth}
            \centering
            \safeincludegraphics[width=\linewidth]{resource_comparison.png}
        \end{column}

        \begin{column}{0.35\textwidth}
            \begin{block}{Đánh đổi}
                \begin{itemize}
                    \item BiLSTM + Attention đạt chất lượng cao nhất nhưng train
                    lâu nhất: 308,47 giây.
                    \item Transformer nhanh nhất: 56,42 giây train và
                    0,186 ms/mẫu.
                    \item Số tham số ba mô hình khá gần nhau do embedding chiếm
                    phần lớn.
                \end{itemize}
            \end{block}
        \end{column}
    \end{columns}
\end{frame}

% =========================================================
\section{4. Kết luận}

\begin{frame}{\textbf{\large Kết luận và hướng phát triển}}
    \begin{columns}[T,onlytextwidth]
        \begin{column}{0.48\textwidth}
            \begin{block}{Kết luận}
                \begin{itemize}
                    \item BiLSTM + Attention là mô hình tốt nhất trong ba kiến
                    trúc đã thử.
                    \item BiLSTM cung cấp ngữ cảnh hai chiều; attention tổng hợp
                    token linh hoạt hơn hidden state cuối.
                    \item Mô hình phù hợp làm kiến trúc chính của project.
                \end{itemize}
            \end{block}
        \end{column}

        \begin{column}{0.48\textwidth}
            \begin{alertblock}{Hạn chế và hướng tiếp theo}
                \begin{itemize}
                    \item Mất cân bằng lớp làm mô hình co cụm về 1 và 5 sao.
                    \item Bật class weights hoặc thử focal loss.
                    \item Chọn checkpoint theo Macro F1.
                    \item Thử ordinal loss và tạo test set độc lập.
                    \item Xuất attention weights để phân tích token.
                \end{itemize}
            \end{alertblock}
        \end{column}
    \end{columns}
\end{frame}

\section{5. Tài liệu tham khảo}

\begin{frame}[allowframebreaks]{\textbf{\large Tài liệu tham khảo}}
    \scriptsize
    \begin{enumerate}
        \item Amazon Reviews Dataset, Kaggle:
        \url{https://www.kaggle.com/datasets/dongrelaxman/amazon-reviews-dataset}.

        \item S. Hochreiter and J. Schmidhuber,
        ``Long Short-Term Memory,'' \textit{Neural Computation}, 1997.

        \item D. Bahdanau, K. Cho, and Y. Bengio,
        ``Neural Machine Translation by Jointly Learning to Align and
        Translate,'' \textit{ICLR}, 2015.

        \item J. Pennington, R. Socher, and C. D. Manning,
        ``GloVe: Global Vectors for Word Representation,''
        \textit{EMNLP}, 2014.

        \item A. Vaswani et al.,
        ``Attention Is All You Need,'' \textit{NeurIPS}, 2017.
    \end{enumerate}
\end{frame}

\begin{frame}[c,plain,noframenumbering]
    \centering
    \Huge\color{hustred}{\textbf{Cảm ơn mọi người\\đã chú ý lắng nghe!}}
\end{frame}

\end{document}
